%document
\documentclass[10pt]{beamer}
%theme
\usetheme{metropolis}
% packages
\usepackage{color}
\usepackage{listings}
\usepackage[main=ngerman, USenglish]{babel}
\usepackage[utf8]{inputenc}
\usepackage{multicol}
\usepackage{csquotes}


% color definitions
\definecolor{mygreen}{rgb}{0,0.6,0}
\definecolor{mygray}{rgb}{0.5,0.5,0.5}
\definecolor{mymauve}{rgb}{0.58,0,0.82}
\definecolor{paleorange}{HTML}{CC7832}

\lstset{
    backgroundcolor=\color{white},
    % choose the background color;
    % you must add \usepackage{color} or \usepackage{xcolor}
    basicstyle=\footnotesize\ttfamily,
    % the size of the fonts that are used for the code
    breakatwhitespace=false,
    % sets if automatic breaks should only happen at whitespace
    breaklines=true,                 % sets automatic line breaking
    captionpos=b,                    % sets the caption-position to bottom
    commentstyle=\color{mygreen},    % comment style
    % deletekeywords={...},
    % if you want to delete keywords from the given language
    extendedchars=true,
    % lets you use non-ASCII characters;
    % for 8-bits encodings only, does not work with UTF-8
    frame=single,                    % adds a frame around the code
    keepspaces=true,
    % keeps spaces in text,
    % useful for keeping indentation of code
    % (possibly needs columns=flexible)
    keywordstyle=\color{blue},       % keyword style
    % morekeywords={*,...},
    % if you want to add more keywords to the set
    numbers=left,
    % where to put the line-numbers; possible values are (none, left, right)
    numbersep=5pt,
    % how far the line-numbers are from the code
    numberstyle=\tiny\color{mygray},
    % the style that is used for the line-numbers
    rulecolor=\color{black},
    % if not set, the frame-color may be changed on line-breaks
    % within not-black text (e.g. comments (green here))
    stepnumber=1,
    % the step between two line-numbers.
    % If it's 1, each line will be numbered
    stringstyle=\color{mymauve},     % string literal style
    tabsize=4,                       % sets default tabsize to 4 spaces
    % show the filename of files included with \lstinputlisting;
    % also try caption instead of title
    language = Java,
	showspaces = false,
	showtabs = false,
	showstringspaces = false,
	escapechar = ,
        morecomment=[s][\textcolor{paleorange}]{@}{\ },
}

\def\ContinueLineNumber{\lstset{firstnumber=last}}
\def\StartLineAt#1{\lstset{firstnumber=#1}}
\let\numberLineAt\StartLineAt



\newcommand{\codeline}[1]{
        \alert{\texttt{#1}}
}

\setmonofont{Fira Code}

% This Document contains the information about this course.

% Authors of the slides
\author{Marcus Köhler}

% Name of the Course
\institute{Java-Kurs}

% Fancy Logo
\titlegraphic{\hfill\includegraphics[height=1.25cm]{../templates/fsr_logo_cropped}}

\usepackage{csquotes}
\usepackage{tikz}
\usepackage{tikz-qtree}

\title{Java}
\subtitle{Scopes, Reference \& JavaDoc}
\date{\today}

\begin{document}

\begin{frame}
\titlepage
\end{frame}

\begin{frame}{Überblick}
    \setbeamertemplate{section in toc}[sections numbered]
    \tableofcontents
\end{frame}

\section{Scopes}

\begin{frame}[fragile]{Scopes}
    \center
    \large{Variablen sind nur innerhalb des Kontexts(d.h. innerhalb der \texttt{\{\}}) verfügbar, in denen sie deklariert wurden. \\
    \pause
    Das ist alles.}
    \pause
    \begin{columns}[T]
        \begin{column}{.5\textwidth}
            Richtig:
            \begin{lstlisting}[gobble=16]
                public class Scope {
                    int a = 5;
                    if(a <= 6) {
                        float f = (float)a;
                        System.out.println(f/1.25f);
                    }
                    System.out.println(a*10);
                }
            \end{lstlisting}
        \end{column}
        \begin{column}{.5\textwidth}
            Falsch:
            \begin{lstlisting}[gobble=16]
                public class Scope {
                    int a = 5;
                    if(a == 6) {
                        float f = (float)a;
                    }
                    System.out.println(f/1.25f);
                    //f is only available inside the if block
                    System.out.println(a*10);
                }
            \end{lstlisting}
        \end{column}
    \end{columns}
\end{frame}

\section{Die Reference}
\subsection{Intro}

\begin{frame}{Die Java-Reference}
    Die Java-Standardbibliothek ist in der sogenannten \textit{Reference} dokumentiert. Jede Java-Version hat ihre eigene Version der Reference;
    die momentan aktuellste Version(Java 10) ist unter \small\url{https://docs.oracle.com/javase/10/docs/api/overview-summary.html} \normalsize zu finden.\\
    \pause
    Die Reference ist hierarchisch gegliedert: \\
    \smallskip \small
    \begin{tikzpicture}
        \Tree   [.Reference 
                    [.\node (level0) {\texttt{java.base}}; 
                        [.\node (level1) {\texttt{java.util}}; 
                            \node (level2) {\texttt{ArrayList}}; {\texttt{HashSet}} {\dots}
                        ] 
                        [.{\texttt{java.lang}}
                            {\texttt{Object}} {\texttt{Math}} {\dots}
                        ]
                        {\dots}
                    ]
                    {\dots}
                ]
            \foreach \Value/\Text in {0/{Module},1/{Package},2/{Klasse}}
            {  
              \node[anchor=west] 
                at ([xshift=-2.5cm,yshift=0.05cm]{level2}|-{level\Value}) 
                {\Text};
            }
    \end{tikzpicture}
\end{frame}

\begin{frame}{Struktur einer Reference-Seite}
    Die Seite einer einzelnen Klasse(oder eines Interfaces) in der Reference ist folgendermaßen strukturiert:
    \begin{enumerate}[<+->]
        \item Eine kurze Zusammenfassung der Klasse, bestehend aus: 
            \begin{itemize}
                \item Der \enquote{Vererbungshierarchie} der Klasse
                \item Gegebenenfalls generische Typvariablen der Klasse
                \item Alle implementierten Interfaces \& alle bekannten erbenden Klassen
                \item Alle Super- und Subinterfaces(nur für Interfaces)
                \item Alle bekannten implementierenden Klassen(nur für Interfaces)
                \item 
            \end{itemize}
    \end{enumerate}
\end{frame}

\end{document}
